% !TeX program = xelatex
\documentclass[UTF8,a4paper,11pt]{ctexart}

\usepackage[margin=2.4cm]{geometry}
\usepackage{booktabs}
\usepackage{tabularx}
\usepackage{longtable}
\usepackage{array}
\usepackage{amsmath}
\usepackage{siunitx}
\usepackage{graphicx}
\usepackage{caption}
\usepackage{xcolor}
\usepackage{enumitem}
\usepackage{hyperref}
\hypersetup{colorlinks=true,linkcolor=black,citecolor=black,urlcolor=blue!60!black}
\captionsetup{font=small,labelsep=quad}
\setlist{itemsep=2pt,topsep=3pt}

\title{\textbf{绝缘卫士：CA-YOLO 绝缘子缺陷检测}\\[0.4em]
\large 真实数据训练、消融、困难工况增广与纯算法验证报告（2026-11-02）}
\author{何周俣\\武汉理工大学 \, 电气工程创新创业实践}
\date{2026 年 11 月 2 日}

\begin{document}
\maketitle

\begin{abstract}
\noindent 本报告记录「基于 CA 注意力机制的绝缘子缺陷检测」这一视觉模块从仿真链路走向真实
数据训练的全过程：完成公开数据集 FBID（1600 张真实巡检照片）的下载、类别映射与划分，
在 RTX 4090 Laptop 上以 YOLOv8n 为基线完成多组 60 epoch 训练，并对三类因素做对照——
(1) 坐标注意力（CoordAtt, CA）的注入位置（不注入 / 注入 SPPF 层 / 注入 backbone 全部 9 层）；
(2) batch size（32 / 64 / 96）；(3) flashover 类别的数据补充（3 倍过采样与同源新增样本合并）。
以 broken 为主验证类别时，CA 的两种注入方式均未表现出可复现的精度增益，差异落在 seed 与
batch 造成的噪声范围之内；batch size 对吞吐影响显著（batch 96 可达 17--19 s/epoch、GPU
利用率 65\%），但对精度并非单调有利。本报告如实记录这些负结果，并给出后续可执行的验证路线。

\vspace{0.6em}
\noindent 在 10-29 版基础上，本版补入\textbf{困难工况增广}与\textbf{输入分辨率}两组实验：
用雾、运动模糊、离焦、暗光、雨痕五种退化合成困难测试集，比较「原始数据（S1）／仅退化数据（S2）／
原始+退化混合（S3）」三种训练配方，并补做 640 与 960 两种输入分辨率的对照。
结论是：\textbf{算法侧（CA、RFA、分辨率）在本数据上都不构成稳定增益，唯一在两个测试条件下
都提升的是数据侧混合增广}（干净集 +1.2\,pt、困难集 +14.1\,pt）；960 分辨率只对清晰图有效
（干净集 +2.8\,pt，困难集 $-4.0$\,pt）。

\vspace{0.6em}
\noindent 在 10-31 版基础上，本版再补入\textbf{不依赖现场硬件的算法侧验证}：
把视觉端的实测误差注入决策闭环（漏检 / 误检）、在调度器里加 K 帧确认与置信度门限、
做 YOLOv8s 容量阶梯与置信度门限扫描、量化相机标定误差对地面落点的影响、
并比较整图与切片推理（NMS / NMM 两种合并）。核心结论有三条：
\textbf{第一}，闭环对漏检不敏感（15\% 漏检下步数不变），对误检极敏感（15\% 误检即不收敛），
决策侧启发式只能减噪，收敛必须靠检测端提高精确率；
\textbf{第二}，相机几何本身把可视带限制在 1.19--3.99\,m，规划容差 0.15\,m 反推要求
俯仰标定优于 $1^\circ$、挂高误差小于 5\,cm；
\textbf{第三}，切片推理只在「目标跨切片」时有收益（跨域整串 +14.6\,pt），域内基本打平。
\end{abstract}

\section{实验目的}

视觉模块的最终任务是：在绝缘子巡检图像中检出 \textbf{normal}（正常）、\textbf{flashover}
（闪络/污秽）与 \textbf{broken}（破损）三类目标，并把缺陷位置交给决策层规划复检机位。

此前版本使用几何投影的仿真检测器（\texttt{src/integration/sim\_detector.py}）打通控制链路，
它只回答「框落在图像哪里」，不回答「识别准不准」。本轮的目标因此是：

\begin{enumerate}[label=(\arabic*)]
  \item 用\textbf{真实标注数据}训练、评测检测器，给出可引用的 mAP 数字；
  \item 验证论文构想中的 \textbf{CA 坐标注意力}是否带来增益，若没有，如实记录；
  \item 找出影响精度的\textbf{非模型因素}（batch、seed、数据分布），为后续实验定口径。
\end{enumerate}

\section{数据集}

\subsection{FBID 数据集}

训练与评测使用 GitHub 公开数据集 \textbf{FBID}（\url{https://github.com/xxlree/FBID}，
EPL-2.0 许可），共 1600 张真实无人机/地面巡检照片。原始标注为 3 类
（\texttt{0 pollution\_flashover}、\texttt{1 broken}、\texttt{2 insulator}），
按项目类别定义重映射为：

\begin{center}
\begin{tabular}{ccc}
\toprule
FBID 原始类别 & 项目类别 & 说明 \\
\midrule
\texttt{2 insulator} & \texttt{0 normal} & 正常绝缘子 \\
\texttt{0 pollution\_flashover} & \texttt{1 flashover} & 闪络/污秽 \\
\texttt{1 broken} & \texttt{2 broken} & 破损 \\
\bottomrule
\end{tabular}
\end{center}

原始 \texttt{train}（1280 张）按「是否含缺陷」分层抽 10\% 作为验证集，原始 \texttt{test}
（320 张）保持独立，最终划分如表~\ref{tab:data}。转换脚本为
\texttt{scripts/prepare\_dataset.py}，产物为 \texttt{训练样本/FBID/yolo/}。

\begin{table}[htbp]
\centering
\caption{FBID 数据划分与标注框统计（按项目类别）}
\label{tab:data}
\begin{tabular}{lccc}
\toprule
split & 图像数 & normal / flashover / broken 框数 & 含缺陷图数 \\
\midrule
train & 1152 & 1331 / 1729 / 823 & 569（flashover）+ 578（broken） \\
val   & 128  & 149 / 212 / 78   & --- \\
test  & 320  & 352 / 507 / 192  & --- \\
\bottomrule
\end{tabular}
\end{table}

\subsection{flashover 补充数据}

flashover 是三类中指标最低的一类（表~\ref{tab:overall} 中 test 仅 0.82），尝试补充数据。
调研发现可用的公开集只有 \texttt{zxc060413/Insulator-defect\_yolo\_format}
（论文 InsuFault-Net，1688 张、4 类）；按文件名比对与 FBID \textbf{82.5\% 同源}，
按图像内容 md5 去重后真正新增仅 \textbf{296 张}（其中 156 张含 flashover）。
因此「外部数据」实际上是同源扩展，不能期望独立的域增益。

\subsection{困难工况数据（人工退化）}

为回答「模型在雾天、雨痕、模糊等恶劣成像下还剩多少精度」，用
\texttt{scripts/make\_degraded.py} 对 FBID 的 train / test 各生成一份退化副本，
五种退化随机叠加（各 1--3 级强度）：大气散射雾、水平运动模糊、离焦模糊、
暗光加噪、雨痕条纹；同一张图随机选 1--2 种，最大边缩至 1280 以控制耗时。
退化图与标签共用同一套 YOLO 标注（框归一化，缩放不影响真值）。

据此定义三种训练配方与两套测试条件：

\begin{center}
\begin{tabular}{lll}
\toprule
名称 & 组成 & 用途 \\
\midrule
S1 & 原始 FBID 1152 张 & 基线 \\
S2 & 仅退化图 1152 张 & 只用困难数据训练 \\
S3 & 原始 + 退化 2304 张 & 混合增广 \\
\midrule
干净测试 & FBID test 320 张 & 正常成像 \\
困难测试 & 同一批 320 张的退化版 & 恶劣成像 \\
\bottomrule
\end{tabular}
\end{center}

三种配方训练出的权重都在两套测试集上评测，且同一配方至少 2 个 seed，
以便把「数据分布」与「随机种子」两种波动分开。

\section{训练环境与超参}

\begin{table}[htbp]
\centering
\caption{训练环境}
\begin{tabular}{ll}
\toprule
项目 & 配置 \\
\midrule
GPU & NVIDIA GeForce RTX 4090 Laptop，16\,GB \\
Python & 3.13.9（独立 venv：\texttt{Demos/.venv}） \\
PyTorch / torchvision & 2.9.1+cu126 / 0.24.1+cu126 \\
Ultralytics & 8.4.192 \\
模型 & YOLOv8n（3.15\,M 参数，CA 注入后 3.17\,M） \\
训练轮数 / 优化器 & 60 epoch / AdamW，余弦退火，warmup 3 \\
batch size & 32 / 64 / 96（对照变量） \\
输入分辨率 & 640 / 960（对照变量） \\
增强 & HSV、水平翻转、mosaic=1.0（Ultralytics 默认组合） \\
数据缓存 & \texttt{cache=ram}（FBID 约 0.6\,GB） \\
随机种子 & seed 0（主要），seed 1（复现） \\
\bottomrule
\end{tabular}
\end{table}

\noindent\textbf{关于 DataLoader 进程数。}受运行环境限制（多进程命名管道被拒绝，
\path{PermissionError: [WinError 5]}），训练与评测统一使用 \texttt{workers=0}，
数据增强在单进程内完成。这也是 batch size 成为吞吐主要杠杆的原因（见第~\ref{sec:speed} 节）。

\section{实验结果}

\subsection{整体与逐类指标（batch 32，两 seed）}

\begin{table}[htbp]
\centering
\caption{整体 mAP（batch 32）。每格为 mAP@0.5 / mAP@0.5:0.95}
\label{tab:overall}
\begin{tabular}{lcccc}
\toprule
run & val mAP50 & val mAP50-95 & test mAP50 & test mAP50-95 \\
\midrule
baseline s0      & 0.9112 & 0.5999 & 0.9056 & 0.6086 \\
baseline s1      & 0.9161 & 0.6069 & 0.9167 & 0.6086 \\
CA@SPPF s0       & 0.9172 & 0.6075 & 0.9031 & 0.6013 \\
CA@SPPF s1       & 0.9124 & 0.6024 & 0.9168 & 0.6089 \\
CA@backbone s0   & 0.9132 & 0.6112 & 0.9062 & 0.6094 \\
\midrule
两 seed 均值（baseline） & 0.9137 & 0.6034 & 0.9112 & 0.6086 \\
两 seed 均值（CA@SPPF）  & 0.9148 & 0.6050 & 0.9100 & 0.6051 \\
\bottomrule
\end{tabular}
\end{table}

CA@SPPF 与基线的两 seed 均值差：val $+0.12$\,pt、test $-0.12$\,pt，而\textbf{同配置
seed 之间的极差为 0.48--0.49\,pt}。也就是说，注意力带来的差异小于随机种子带来的波动，
在该数据与超参下\textbf{测不出 CA 的正向作用}。

\begin{table}[htbp]
\centering
\caption{逐类 mAP@0.5（batch 32，seed 0）}
\begin{tabular}{lcccccc}
\toprule
run & \multicolumn{2}{c}{normal} & \multicolumn{2}{c}{flashover} & \multicolumn{2}{c}{broken} \\
\cmidrule(lr){2-3}\cmidrule(lr){4-5}\cmidrule(lr){6-7}
 & val & test & val & test & val & test \\
\midrule
baseline        & 0.995 & 0.995 & 0.800 & 0.820 & 0.938 & 0.902 \\
CA@SPPF         & 0.995 & 0.995 & 0.804 & 0.813 & 0.953 & 0.901 \\
CA@backbone     & 0.995 & 0.995 & 0.804 & 0.818 & 0.941 & 0.905 \\
\bottomrule
\end{tabular}
\end{table}

normal 已达到饱和（0.995），瓶颈集中在 flashover（0.80--0.82）与 broken（0.90--0.95）。
其中 flashover 与「污秽」在视觉上高度混淆、标注本身也不稳定，更适合作为数据质量问题记录，
而不适合作为算法验证目标；因此后续以 \textbf{broken}（结构破损，标注可信）作为主验证类别。

\subsection{broken 主指标：batch 档位 $\times$ CA 注入位置}

\begin{table}[htbp]
\centering
\caption{broken 单类指标。每格为 mAP@0.5 / mAP@0.5:0.95，同一 batch 内可直接比较}
\label{tab:broken}
\begin{tabular}{llcccc}
\toprule
batch & run & val mAP50 & val mAP50-95 & test mAP50 & test mAP50-95 \\
\midrule
32 & baseline      & 0.9380 & 0.5223 & 0.9015 & 0.5018 \\
32 & CA@SPPF       & 0.9525 & 0.5314 & 0.9013 & 0.4928 \\
32 & CA@backbone   & 0.9408 & \textbf{0.5519} & 0.9050 & 0.5108 \\
\midrule
64 & baseline      & 0.9446 & 0.5285 & 0.9081 & 0.5048 \\
64 & CA@SPPF       & 0.9440 & 0.5044 & 0.9048 & 0.5014 \\
64 & CA@backbone   & 0.9343 & 0.5038 & 0.9087 & 0.5104 \\
\midrule
96 & baseline      & \textbf{0.9514} & 0.5233 & \textbf{0.9095} & 0.5021 \\
96 & CA@backbone   & 0.9167 & 0.4988 & 0.9024 & 0.4808 \\
\bottomrule
\end{tabular}
\end{table}

这是本报告最重要的负面结果：

\begin{enumerate}[label=(\arabic*)]
  \item \textbf{CA 注入位置没有可复现的增益。} batch 32 上 CA@backbone 的
        broken mAP50-95 比基线高 3.0\,pt（val），但同配置在 batch 64 变为 $-2.5$\,pt、
        batch 96 变为 $-2.4$\,pt；batch 32 上 CA@SPPF 的 mAP50 高 1.45\,pt，
        到 batch 64 反低 0.06\,pt。差异没有一致方向。
  \item \textbf{差异量级小于噪声。} 同 batch 内三档 broken mAP50 的最大差距为 1.5\,pt；
        而 seed 之间为 0.9--2.1\,pt、batch 之间为 0.7--1.2\,pt。1.5\,pt 落在噪声带内，
        不能据此宣称「某个注入位置更好」。
  \item \textbf{batch 对精度不是单调有利。} broken mAP50 随 batch 上升
        （0.9380 $\to$ 0.9446 $\to$ 0.9514），但 test 几乎不动
        （0.9015 $\to$ 0.9081 $\to$ 0.9095）；整体 mAP50-95 反而是 batch 64 最好。
\end{enumerate}

\subsection{吞吐：batch 96 的真实收益}
\label{sec:speed}

\begin{table}[htbp]
\centering
\caption{训练吞吐（640 分辨率、60 epoch、单进程数据管线）}
\begin{tabular}{lccc}
\toprule
batch & 每轮耗时 & 显存占用 & GPU 利用率 \\
\midrule
32 & $\sim$30--35\,s & 5.4\,GB & $\sim$40\,\% \\
64 & $\sim$30\,s      & 7.9\,GB & $\sim$45\,\% \\
96 & \textbf{17--19\,s} & 11.9\,GB & \textbf{$\sim$65\,\%} \\
\bottomrule
\end{tabular}
\end{table}

batch 96 把一轮 60 epoch 从约 30 分钟压到约 15 分钟，是本轮唯一稳定可复现的「提升」；
但它在精度上并不占优。若要继续提速，只能改用多进程数据加载（本环境不支持）或
减少增强开销，而不是继续加大 batch。

\subsection{flashover 补充数据实验}

\begin{table}[htbp]
\centering
\caption{flashover 补充实验（每格为整体 mAP@0.5 / flashover mAP@0.5 / broken mAP@0.5）}
\label{tab:flash}
\begin{tabular}{llcccc}
\toprule
split & run & batch & mAP50 & flashover & broken \\
\midrule
val  & baseline            & 32 & 0.9112 & 0.800 & 0.938 \\
val  & E2 flashover$\times$3 & 32 & 0.9193 & \textbf{0.852} & 0.911 \\
val  & baseline            & 64 & \textbf{0.9234} & 0.830 & 0.945 \\
val  & E3 MERGED+$\times$3   & 64 & 0.9149 & \textbf{0.853} & 0.897 \\
val  & E4 CA@backbone      & 64 & 0.9161 & 0.819 & 0.934 \\
\midrule
test & baseline            & 32 & 0.9056 & 0.820 & 0.902 \\
test & E2 flashover$\times$3 & 32 & 0.9043 & 0.830 & 0.888 \\
test & baseline            & 64 & 0.9125 & 0.834 & 0.908 \\
test & E3 MERGED+$\times$3   & 64 & \textbf{0.9199} & \textbf{0.856} & 0.909 \\
test & E4 CA@backbone      & 64 & 0.9027 & 0.804 & 0.909 \\
\bottomrule
\end{tabular}
\end{table}

补充 flashover 数据确实把该单类在两个划分上都推高约 2.2\,pt（E3-b64），test 整体 +0.74\,pt；
代价是 val 的 broken 掉 4.8\,pt——过采样挤占了其他类别的采样权重，整体基本持平。
这从另一个侧面说明：\textbf{当前瓶颈是数据分布，不是注意力结构}。

\subsection{困难工况增广与分辨率：完整矩阵}

\begin{table}[htbp]
\centering
\caption{六组配方 $\times$ 两种测试条件。每格为 mAP@0.5 / mAP@0.5:0.95，逐类列为 mAP@0.5}
\label{tab:matrix}
\setlength{\tabcolsep}{4pt}
\begin{tabular}{lcccccc}
\toprule
配方 & 分辨率 & 干净·全体 & 困难·全体 & 干净·broken & 困难·broken \\
\midrule
S1（3 seed 均值）   & 640 & 0.910 / 0.603 & 0.746 / 0.480 & 0.902 & 0.674 \\
S2（2 seed 均值）   & 640 & 0.890 / 0.581 & 0.865 / 0.557 & 0.899 & 0.868 \\
S3（2 seed 均值）   & 640 & \textbf{0.922 / 0.606} & 0.887 / 0.579 & \textbf{0.927} & \textbf{0.897} \\
S3 + RFA（1 seed）  & 640 & 0.913 / 0.602 & \textbf{0.890 / 0.582} & 0.922 & 0.904 \\
\midrule
S1                  & 960 & \textbf{0.938 / 0.630} & 0.706 / 0.452 & \textbf{0.932} & 0.632 \\
S3                  & 960 & 0.936 / 0.619 & 0.892 / 0.575 & 0.916 & 0.871 \\
\bottomrule
\end{tabular}
\end{table}

\begin{enumerate}[label=(\arabic*)]
  \item \textbf{混合增广（S3）在两个测试条件下同时变好。} 640 分辨率下，干净集 $+1.2$\,pt
        （0.910$\to$0.922）、困难集 $+14.1$\,pt（0.746$\to$0.887），broken 困难集
        $+22.3$\,pt。单 seed 阶段曾观察到「S3 干净集 $-3$\,pt」，补齐 2--3 个 seed 后
        该差异消失——它就是 seed 噪声（S1 三个 seed 自身跨越 0.904--0.915）。
  \item \textbf{仅用退化数据（S2）是「拿干净集换困难集」。} 干净集 $-2.0$\,pt、
        flashover $-5.7$\,pt，困难集 $+11.9$\,pt。若部署场景确定是恶劣天气，S2 可用；
        否则 S3 更稳。
  \item \textbf{分辨率只在清晰图上换精度。} S1 从 640 提到 960，干净集 $+2.8$\,pt
        （0.910$\to$0.938）、broken $+3.0$\,pt；但困难集 $-4.0$\,pt（0.746$\to$0.706）、
        broken $-4.2$\,pt。原因是人工退化（雾、模糊）本身压掉了高频细节，
        高分辨率的优势在输入侧就被抹平，同时放大了训练与测试的分布差。
  \item \textbf{RFA 叠在 S3 上没有增益。} 相对 S3 均值，干净 $-0.9$\,pt、困难 $+0.4$\,pt、
        broken $-0.5$ / $+0.8$\,pt，全部落在 seed 极差（0.9--2.2\,pt）之内；
        结合前面 CA 的结果，\textbf{结构侧在本数据上没有可复现收益}。
  \item \textbf{部署选型。} 晴天/清晰成像用 S1@960（干净 0.938、broken 0.932）；
        恶劣天气用 S3@640（困难 0.887、broken 0.897）。两者困难集相差 19\,pt，
        说明「一套权重走天下」在这个任务上不成立，应按场景选型或做双模型切换。
\end{enumerate}

\subsection{视觉误差如何传到决策线}

识别精度只决定「看见多少」，闭环能不能跑完取决于漏检与误检怎么进调度器。用
\texttt{scripts/sim\_noise.py} 按实测 P/R 给仿真检测器注入两类错误，其余链路不变：

\begin{table}[htbp]
  \centering
  \caption{漏检 / 误检对闭环的影响（完美检测基线：789 步、7 任务、0 碰撞）}
  \label{tab:noise}
  \begin{tabular}{lcccc}
    \toprule
    注入 & 步数 & 建任务 & 完成 & 碰撞 \\
    \midrule
    15\% 漏检 & 789 & 7 & 7 & 0 \\
    30\% 漏检 & 789 & 7 & 7 & 0 \\
    15\% 误检 & 3000（不收敛） & 44 & 6 & 0 \\
    30\% 误检 & 3000（不收敛） & 48 & 13 & 11 \\
    \bottomrule
  \end{tabular}
\end{table}

漏检几乎无影响——多帧观测与缺陷点去重把偶发漏检吃掉了；误检是杀手：15\% 误检就能把
任务队列灌满、巡检任务被饿死。这条结果修正了「识别越准越好」的直觉判断：
\textbf{对闭环而言，误检的边际代价远高于漏检}。

\subsection{K 帧确认与置信度门限的实际效果}

针对上节，在调度器里加了「同一缺陷点要在 $K$ 个不同帧里被看到才建任务」的确认机制，
并与缺陷置信度门限组合成 $3\times2$ 的网格。

\begin{table}[htbp]
  \centering
  \caption{K 帧确认：建任务数 / 幽灵任务数 / 完成任务数（单 seed）}
  \label{tab:kframe}
  \begin{tabular}{ccc}
    \toprule
    $K$ & 15\% 误检 & 30\% 误检 \\
    \midrule
    1 & 40 / 37 / 6 & 46 / 43 / 7 \\
    2 & 36 / 33 / 13 & 45 / 42 / 14 \\
    3 & 36 / 33 / 19 & 38 / 35 / 19 \\
    \bottomrule
  \end{tabular}
\end{table}

$K=2$ 是性价比点：幽灵任务 $-11\%$，真缺陷几乎不丢；$K=3$ 再多压 8\%，但完成数停滞。
门限则更糟：0.85 配 $K=2$ 时确实零幽灵，但 3 个真缺陷只复检到 1 个——因为真缺陷的
置信度本来就在 0.6--0.9 区间。结论必须如实写：\textbf{决策侧的启发式只能减噪，
不能把闭环拉回收敛；误检要在检测端抑制}（更高 P、更强 NMS、由检测器输出跨帧轨迹）。

\subsection{容量阶梯与置信度门限}

模型侧再试一次「换大模型」：S3 配方把 YOLOv8n 换成 YOLOv8s（60\,epoch、batch 48）。

\begin{table}[htbp]
  \centering
  \small
  \setlength{\tabcolsep}{4pt}
  \caption{容量阶梯（S3 配方，640）}
  \label{tab:capacity}
  \begin{tabular}{lcccc}
    \toprule
    配方 & 干净 .5 / .5:.95 & 困难 .5 / .5:.95 & broken 干净 & broken 困难 \\
    \midrule
    S3 + YOLOv8n（2 seed 均值） & 0.922 / 0.606 & 0.887 / 0.579 & 0.927 & 0.897 \\
    S3 + YOLOv8s（1 seed）      & 0.911 / 0.612 & 0.888 / 0.586 & 0.927 & 0.925 \\
    \bottomrule
  \end{tabular}
\end{table}

干净集 $-1.1$\,pt、困难集 $+0.1$\,pt，只有 mAP@.5:.95 涨 0.6--0.7\,pt。
2304 张图撑不起更大容量——这与前面 CA/RFA 的结论一致：\textbf{瓶颈不在模型容量，
而在数据分布}。

置信度门限的扫描（\texttt{sweep\_broken\_conf.py}）给出另一条实用结论：
F1 在 0.05--0.45 是一条平台（0.895--0.899，最优 0.30 时 P\,0.962 / R\,0.844），
0.55 起塌陷，0.85 时 recall 只剩 0.031。部署把门限留在 \textbf{0.25--0.30} 即可，
抬门限只会换漏检——这也说明上一节仿真的 15--30\% 误检属于压力测试，
实测 broken 的精确率已经到 0.95。

\subsection{相机标定鲁棒性：几何先定死一条带}

检测框只有落到地面坐标系才有意义。用 \texttt{scripts/calib\_robust.py} 对相机内外参
做扰动扫描（19\,481 条射线，全 numpy 向量化）：

\begin{table}[htbp]
  \centering
  \caption{相机参数扰动 $\rightarrow$ 落点误差 p95（挂高 1.5\,m、俯仰 25$^\circ$）}
  \label{tab:calib}
  \begin{tabular}{lcc}
    \toprule
    扰动 & 全体 p95 & 远端（3.3--4.0\,m）p95 \\
    \midrule
    俯仰 $+1^\circ$ & 0.14\,m & 0.16\,m \\
    俯仰 $+2^\circ$ & 0.26\,m & 0.30\,m \\
    挂高 $+10$\,cm & 0.22\,m & 0.26\,m \\
    横滚 $+3^\circ$ & 0.21\,m & 0.32\,m \\
    $f_x +2\%$ & 0.03\,m & 0.04\,m \\
    检测框中心 $+5$\,px & 0.04\,m & 0.04\,m \\
    组合（俯 $+2^\circ$、高 $+0.1$\,m、滚 $+1^\circ$） & 0.09\,m & 0.16\,m \\
    \bottomrule
  \end{tabular}
\end{table}

两个容易被忽略的事实：其一，这套挂高与俯仰下地板可见带只有 \textbf{1.19--3.99\,m}
（光轴落点 3.22\,m），仿真里 \texttt{max\_range\_m}$=6$ 有一半用不上，4\,m 之外的
检测框根本投不到地面；其二，规划层 \texttt{goal\_tolerance} 是 0.15\,m，
反推得到\textbf{俯仰标定要优于 $1^\circ$、挂高误差要小于 5\,cm}。
组合扰动反而更小，是因为俯仰把落点拉近、挂高把它推远，两项会互相抵消——
实测标定不能只校准一个量。

\subsection{切片推理与 NMM：只在「目标跨切片」时有效}

大图直接缩到 640 会抹平小缺陷，切片推理是一种补救。两种合并方式差别很大：
NMS 重叠就抑制，会把被切开的同一根绝缘子丢一半；NMM（Non-Maximum Merging）
把重叠/包含的同类框并成外接框，能把切开的两半拼回去。

\begin{table}[htbp]
  \centering
  \caption{recall@0.5：整图 vs 2$\times$2 切片（overlap 0.35，imgsz 640）}
  \label{tab:tiled}
  \begin{tabular}{lcccc}
    \toprule
    数据集 / 类别 & 整图 & 切片 + NMS & 切片 + NMM & GT 数 \\
    \midrule
    FBID 全体 & 88.24\% & 89.46\% & 89.46\% & 408 \\
    \quad normal & 100.0\% & 98.5\% & 99.2\% & 132 \\
    \quad flashover & 77.8\% & 82.8\% & 82.3\% & 203 \\
    \quad broken & 95.9\% & 91.8\% & 91.8\% & 73 \\
    \midrule
    CPLID 整串 全体 & 24.27\% & 28.87\% & \textbf{38.91\%} & 234 \\
    \quad normal & 30.2\% & 35.9\% & \textbf{48.4\%} & 192 \\
    \quad broken & 0.0\% & 0.0\% & 0.0\% & 47 \\
    \bottomrule
  \end{tabular}
\end{table}

域内 FBID 上 NMM 与 NMS 打平（绝缘子只占画面一小块，很少跨切片边界）；
跨域的 CPLID 整串图才见真章：NMM 比 NMS 高 10\,pt、比整图高 14.6\,pt。
但 broken 三档全 0——模型在该域根本认不出破损，\textbf{切片只算止血，不解决域差}。

\begin{quote}
\small 值得单独记一笔：这一组数字第一次跑出来时 NMM 明显更差，排查发现是
\texttt{iou\_matrix} 把 \texttt{[cls, cx, cy, w, h, conf]} 当成 \texttt{cxcywh} 在算，
等于拿类别号当坐标。指标函数必须先做「同框自匹配 IoU $=1$」的自检，
否则整套对照都是无效的。
\end{quote}

\section{结论与口径}

\begin{enumerate}[label=(\arabic*)]
  \item 真实数据上，YOLOv8n 在 FBID 的整体 test mAP@0.5 为 \textbf{0.906--0.917}（640）、
        \textbf{0.936--0.938}（960），broken 单类为 \textbf{0.90--0.95}；
        在人工退化测试集上则降到 \textbf{0.706--0.892}，取决于训练配方。这是可以对外引用的数字。
  \item \textbf{算法侧没有可复现的稳定增益}：CA 的两种注入位置（$\pm0.12$\,pt）、
        RFA（$\le0.9$\,pt）、960 分辨率（干净 $+2.8$ / 困难 $-4.0$\,pt）都落在 seed 极差
        附近或只在单一条件下成立。论文与答辩中应把 CA/RFA 表述为「结构上的可选增强」，
        而不是「精度提升点」。
  \item \textbf{数据侧有可复现的增益}：原始+退化混合增广在两个测试条件下同时提升
        （干净 $+1.2$、困难 $+14.1$\,pt），是唯一「全天候」的有效的配置；
        仅退化数据的配方则会牺牲干净集精度。这一条比 10-29 版「数据侧远大于结构侧」
        的说法更精确——混合增广不是零和，它两个都赚。
  \item batch size 与 seed 的交互大于 CA/RFA 的影响；cross-batch 的数字\textbf{不可}
        直接比较，所有对照必须同 batch、同 seed 数、同数据划分。
  \item flashover 的短板来自标注与视觉歧义（困难集 0.63--0.77），补同源公开数据只能缓解、
        不能根治；真正的下一步是现场巡检图，以及按场景选型（或双模型切换）。
  \item \textbf{闭环对漏检不敏感、对误检极敏感}：15\% 漏检下步数仍为 789（与完美检测相同），
        15\% 误检就要 3000 步仍不收敛。调度侧的 K 帧确认（$K=2$）能把幽灵任务压 11\%、
        并把 30\% 误检下的碰撞从 11 次降到 0，但{\bf 没有一档能收回到 789 步}——
        误检必须在检测端抑制。
  \item \textbf{容量与门限都已封顶}：YOLOv8s 相比 v8n 干净集 $-1.1$\,pt、困难集 $+0.1$\,pt，
        2304 张图撑不起更大模型；置信度门限的 F1 在 0.05--0.45 是平台（最优 0.30），
        部署把门限留在 0.25--0.30，抬门限只换漏检。
  \item \textbf{相机几何限制了整条链路}：挂高 1.5\,m + 俯仰 25$^\circ$ 时地板可见带只有
        1.19--3.99\,m，4\,m 外的检测框投不到地面；规划容差 0.15\,m 反推要求俯仰标定
        $\le1^\circ$、挂高误差 $\le5$\,cm，两项偏差会互相抵消，实测不能只校准一个。
  \item \textbf{切片推理是条件性增益}：域内 FBID 上 NMM 与 NMS 打平（89.5\%），
        跨域整串 CPLID 上 NMM 38.9\% $>$ NMS 28.9\% $>$ 整图 24.3\%；
        但 broken 三档全 0，说明\textbf{切片只算止血，不解决域差}。
        附带教训：指标函数必须先用「同框自匹配 IoU $=1$」自检，否则整套对照无效。
\end{enumerate}

\section{复现与产物}

\noindent 以下命令均在 \texttt{Demos/insulator\_guardian} 下执行；\texttt{python} 指训练虚拟环境
\texttt{..\textbackslash.venv\textbackslash Scripts\textbackslash python.exe}（依赖装在该 venv 内，
系统 Python 未装 torch / ultralytics）。

\begin{verbatim}
# 数据准备（FBID 原始标注 -> 项目 3 类 YOLO 划分）
python scripts/prepare_dataset.py --src ../训练样本/FBID/raw \
       --dst ../训练样本/FBID/yolo

# 训练：基线 / CA@SPPF / CA@backbone（batch 96 示例）
python scripts/train_vision.py --run-name real_baseline_b96_v8n \
       --no-ca --epochs 60 --batch 96 --cache ram --seed 0
python scripts/train_vision.py --run-name real_ca_sppf_b64_v8n \
       --epochs 60 --batch 64 --cache ram --seed 0
python scripts/train_vision.py --run-name real_ca_backbone_b96_v8n \
       --epochs 60 --batch 96 --cache ram --seed 0 --ca-inject backbone

# 评测（val / test）
python scripts/evaluate_vision.py --weights ../results/runs/<run>/weights/best.pt \
       --split test --device 0

# 汇总表
python scripts/summarize_runs.py              # training_summary.md
python scripts/summarize_runs.py --focus broken  # broken_summary.md

# 困难工况数据与实验矩阵（本版新增）
python scripts/make_degraded.py --dataset both     # -> 训练样本/DEG
python scripts/train_vision.py --run-name real_deg_v8n \
       --data-yaml ../训练样本/DEG/data_s2.yaml \
       --no-ca --epochs 60 --batch 64 --cache ram --seed 0
python scripts/train_vision.py --run-name real_degmix_b64_s1 \
       --data-yaml ../训练样本/DEG/data_s3.yaml \
       --no-ca --epochs 60 --batch 64 --cache ram --seed 1 \
       --imgsz 640

# 逐类 P/R 对照（多权重 × 多测试集）
python scripts/eval_pr_table.py --device 0 \
       --weights "S1=../results/runs/real_baseline_b64_v8n/weights/best.pt" \
       --weights "S3=../results/runs/real_degmix_b64_s1/weights/best.pt" \
       --test "干净=../训练样本/FBID/yolo/data.yaml" \
       --test "困难=../训练样本/DEG/data_s1_testdeg.yaml"
\end{verbatim}

\noindent 产物位置：
\begin{itemize}
  \item 权重与训练曲线：\texttt{Demos/results/runs/<run>/}；
  \item 评测 JSON：\texttt{Demos/results/benchmarks/vision\_eval\_*.json}；
  \item 汇总表：\texttt{training\_summary.md}、\texttt{broken\_summary.md}；
  \item 逐类 P/R：\texttt{pr\_table\_2026-10-31.md} 与 \texttt{pr960/}；
  \item 工作日志（含全部实验矩阵与踩坑记录）：\texttt{Demos/WORKLOG.md}。
\end{itemize}

\section{已知限制}

\begin{itemize}
  \item 数据集为公开集，与决赛现场工况、拍摄角度不完全一致；同源数据已接近天花板。
  \item 困难测试集是\textbf{人工合成退化}（雾/运动模糊/离焦/暗光/雨痕），只代表一种退化
        分布，不等于真实雾天、雨天；真实数据的退化更复杂，且伴随光照与运动变化。
  \item S2/S3 为 2 seed、S3+RFA 与两档 960 各 1 seed，$\pm1$\,pt 级别的差异仍需更多
        seed 才能下更硬的结论；S1 已到 3 seed。
  \item 受环境限制，数据加载为单进程，吞吐不代表硬件上限。
  \item conf 阈值扫描已在困难集上跑完（最优 0.30、F1 0.899），但仍只用单 seed、单权重；
        换权重后最优点可能小幅漂移。
  \item 本轮全部验证是\textbf{算法侧与仿真闭环}，没有实物、没有现场数据：相机标定的结论来自
        几何仿真，切片与 NMM 的跨域结论来自 CPLID 的 143 张图，都还不能替代真机实验。
  \item 误检抑制只在调度层做了 K 帧确认与门限，检测端（更强 NMS、跨帧轨迹、难负样本挖掘）
        尚\textbf{未}实现，这是闭环收敛的真正瓶颈。
\end{itemize}

\clearpage
\appendix
\section{创新点与论证缺口}

本附录把项目对外主张的每条创新点，与截至本轮的实验证据逐条对照，标出证据不足、口径
风险与处置建议。它不改变正文的实验结论，只回答一个问题：\textbf{现在能不能写}. 各条
对应的原始数字见 \texttt{Demos/results/benchmarks/} 与 \texttt{Demos/WORKLOG.md}。

\begin{scriptsize}
\begin{longtable}{@{}>{\raggedright\arraybackslash}p{2.15cm}
                     >{\raggedright\arraybackslash}p{3.95cm}
                     >{\raggedright\arraybackslash}p{4.05cm}
                     >{\raggedright\arraybackslash}p{3.45cm}@{}}
\caption{各创新点的现有证据、问题与处置}\label{tab:novelty}\\
\toprule
创新点 & 现有证据 & 问题 & 处置建议 \\
\midrule
\endfirsthead
\multicolumn{4}{@{}l}{\footnotesize 表~\thetable（续）}\\[2pt]
\toprule
创新点 & 现有证据 & 问题 & 处置建议 \\
\midrule
\endhead
\bottomrule
\endfoot
CA 坐标注意力
  & 3 档 batch × 3 种注入共 10 组 run；同 batch 内差 $\le 1.5$\,pt，跨 batch 方向相反
  & 无稳定增益；CA 本身是 2021 年已有方法，非新结构
  & 只写「可选结构增强」；或补 SE / CBAM / ECA 对照，每档 $\ge 3$ seed \\
RFA 多尺度聚合
  & 已接进 YOLOv8n 主干并用 3 种数据配方训练：单独用 $+0.2$ / $+1.8$\,pt，
    叠在 S3 上相对 S3 均值 $-0.9$ / $+0.4$\,pt，均在 seed 极差内
  & 有实验、无增益；5 层融合带来的参数量与延迟都得付
  & 从创新点中删除，或作为「可切换模块」写成工程选项 \\
困难工况增广（S1/S2/S3）
  & 2--3 seed：S3 干净 $+1.2$ / 困难 $+14.1$\,pt，broken 困难 $+22.3$\,pt
  & 困难测试集是人工合成退化，不等于真实雾天雨天
  & 可作为本轮唯一站得住的正向结论，配真实现场图验证 \\
频域增强（CWT / VMD）
  & 真实裁剪 32$\times$32，空间特征 40.78\% $\rightarrow$ 融合 38.98\%
  & 负增益
  & 定位为光照/退化归一化分支，不与检测 mAP 挂钩 \\
图结构 / 小样本（MSA-GCN、cDMTH）
  & 3 类手工特征 45.7\% / 42.6--45.9\%，随机基线 33\%
  & 只证明流程可运行，未接入主干，不构成检测增益
  & 降为探索性实验，不与检测指标并列 \\
cDCGAN 数据增广
  & 60 / 300 epoch，生成样本仍为低频色块
  & 生成质量未收敛，不能作为增广数据使用
  & 明写「管线跑通、质量不可用」 \\
缺陷驱动复检调度 + A*/DWA 联动
  & 仿真闭环 789 步 / 7 任务 / 0 碰撞；随机图基准 50 例成功率 94\%，2 例碰撞
  & 组件均为成熟算法，新意只在联动；证据全部来自仿真
  & \textbf{建议作为创新主线}，并补半实物或真机验证 \\
数据集贡献
  & 使用公开集 FBID，仅做类别映射与重新划分
  & 无自有数据、无新增标注
  & 不作为数据集贡献主张 \\
误差传播与调度加固
  & 15\% 误检即不收敛；$K=2$ 把幽灵任务压 11\%、30\% 误检下碰撞 11$\rightarrow$0，但收不回 789 步
  & 说明闭环瓶颈在检测端精确率，决策侧启发式只是减噪
  & 写成「系统层分析」而非创新点，避免把调度启发式当卖点 \\
相机标定鲁棒性
  & 可视带 1.19--3.99\,m；0.15\,m 容差反推俯仰 $\le1^\circ$、挂高 $\le5$\,cm
  & 全部来自几何仿真，未做实物标定
  & 作为工程约束写进建车/装机要求，不算算法创新 \\
切片推理与 NMM
  & 域内打平；跨域整串 NMM 38.9\% 比整图 24.3\% 高 14.6\,pt，broken 仍 0
  & 只对超大目标有效，且在跨域场景仍不可用
  & 作为「大图补救」写进工程选项，不写成精度提升 \\
\end{longtable}
\end{scriptsize}

\noindent 另有四条跨条目的共性缺口，任何一条创新点想站住都必须先补：

\begin{enumerate}[label=(\arabic*)]
  \item 主实验 seed 数不齐（S1 已 3 个，S2/S3 为 2 个，960 与 S3+RFA 各 1 个），
        缺显著性检验；$\pm1$\,pt 量级的差异不足以下结论。
  \item 无 SOTA 对比：未与 YOLOv8s/m、YOLOv10/11、RT-DETR 等横向比较，
        也未在 CPLID / IDID 等公开集上交叉验证。
  \item test 仅 320 图（broken 192 框），val/test 之间 0.5--1.5\,pt 的差异属抽样抖动。
  \item 困难测试集是人工合成退化，真实雾天/雨天还带光照与运动变化，域差距未覆盖。
\end{enumerate}

\enlargethispage{4\baselineskip}
\noindent 综上，本轮真正立得住的有两层：\textbf{系统层}（检测 $\rightarrow$ 世界坐标
$\rightarrow$ 复检机位 $\rightarrow$ 路径规划 $\rightarrow$ 任务完成的闭环，仿真闭环 789 步 /
7 任务 / 0 碰撞）与\textbf{数据层}（混合增广在干净与困难条件下同时提升）。
建议创新点叙事以这两条为主线，CA / RFA / 分辨率降为设计选择而非卖点。

\end{document}
